% Documentación de la Etapa 2
% Compilar con:  pdflatex Informe_etapa2.tex
\documentclass[11pt,a4paper]{article}

\usepackage[utf8]{inputenc}
\usepackage[margin=2.5cm]{geometry}
\usepackage{url}

\title{Etapa 2 - Analizador Sintáctico de MiniJava}
\author{Compiladores e Intérpretes --- UNS, 2026}
\date{}

% Esta instalacion TeX no tiene las fuentes TS1 (text companion): se evitan los
% simbolos que las requieren (vinetas de itemize, llaves via \{ en \texttt).
\renewcommand{\labelitemi}{$\bullet$}

\begin{document}

\sloppy

\maketitle

\noindent
Implementación de un analizador sintáctico descendente recursivo para MiniJava, cliente del
analizador léxico desarrollado en la Etapa 1. La gramática original de MiniJava no es LL(1),
por lo que primero se la transformó (eliminación de recursión a izquierda y factorización
simple y profunda) obteniendo una gramática equivalente LL(1), sobre la que se aplicó la
estrategia recursiva simple: un método por cada no terminal, selección de la producción
según los primeros de sus partes derechas y consumo de terminales mediante \texttt{match}.

\section{Esquema general}

El módulo principal recibe el archivo fuente y se lo asocia al manejador de archivos y este
al léxico (como en la Etapa 1). Luego crea al sintáctico con la dependencia al léxico y le
pide que comience el análisis. Si el análisis finaliza sin errores, informa el mensaje de
éxito; si se produce un error sintáctico o léxico, el sintáctico termina propagando una
excepción que el módulo principal captura y reporta. A diferencia de la Etapa 1, el
principal no imprime la secuencia de tokens: su interfaz se limita a disparar el análisis y
reportar el resultado.

\section{Estructura del proyecto}

{\footnotesize
\begin{verbatim}
codigo/
|-- Compilar.sh / Compilar.bat          Compila y genera Compilador.jar
|-- CorrerTests.sh / CorrerTests.bat    Compila y ejecuta los testers JUnit
|-- lib/                                Jars de JUnit 4 y Hamcrest (se autodescargan)
|-- src/main/java/
|   |-- moduloprincipal/
|   |   \-- ModuloPrincipal.java        Modulo principal (interfaz con el usuario)
|   |-- analizadorlexico/
|   |   |-- AnalizadorLexico.java       Automata finito determinista (Etapa 1, sin cambios)
|   |   |-- Token.java                  Representacion de un token (nombre, lexema, linea)
|   |   \-- ExcepcionLexica.java        Excepcion propagada ante un error lexico
|   |-- analizadorsintactico/
|   |   |-- AnalizadorSintactico.java   Parser descendente recursivo (un metodo por NT)
|   |   \-- ExcepcionSintactica.java    Excepcion propagada ante un error sintactico
|   \-- sourcemanager/
|       |-- SourceManager.java          Interfaz provista por la catedra
|       |-- SourceManagerImpl.java      Implementacion provista (por lineas, referencia)
|       \-- SourceManagerEficiente.java Implementacion propia por caracteres (se usa)
|-- src/test/java/test/                 Testers JUnit provistos por la catedra
\-- resources/
    |-- sinErrores/                     Casos de prueba sin errores
    \-- conErrores/                     Casos de prueba con errores
\end{verbatim}
}

\section{Componentes}

\subsection{Módulo principal (ModuloPrincipal)}

Recibe por parámetro la ruta del archivo fuente MiniJava (acepta cualquier extensión):

\begin{verbatim}
java -jar Compilador.jar programa1.java
\end{verbatim}

Asocia el archivo a un \texttt{SourceManagerEficiente}, crea el \texttt{AnalizadorLexico},
crea el \texttt{AnalizadorSintactico} con esa dependencia y le solicita el análisis. Al
finalizar sin errores imprime el mensaje de éxito y \texttt{[SinErrores]}. Ante un error
sintáctico o léxico captura la excepción correspondiente, reporta el mensaje y el código de
error y finaliza la ejecución, como exigen las pautas de esta etapa.

\subsection{Analizador sintáctico (AnalizadorSintactico)}

Implementa la estrategia recursiva simple sobre la gramática LL(1) de la sección
\ref{sec:gramatica}:

\begin{itemize}
  \item \textbf{Un método por cada no terminal} de la gramática. El cuerpo codifica las
        partes derechas de sus producciones: por cada terminal llama a \texttt{match}, y
        por cada no terminal llama a su método.
  \item \textbf{Selección por primeros}: cuando un NT tiene varias producciones, se elige
        la primera cuya parte derecha tiene al token actual en su conjunto de primeros
        (conjuntos precalculados como constantes, p.\,ej.\ los primeros de un tipo o de una
        expresión). Si el NT tiene una producción $\epsilon$, el esquema simple no hace
        nada en la rama \texttt{else}.
  \item \textbf{match}: compara el token actual con el terminal esperado; si coinciden pide
        el próximo token al léxico, y si no lanza \texttt{ExcepcionSintactica} con el
        lexema y la línea del token actual.
  \item \textbf{Arranque}: el método \texttt{analizar()} carga el primer token del léxico y
        llama al método del símbolo inicial \texttt{<Inicial>}.
  \item Un mapa de descripciones de tokens permite generar mensajes del estilo de las
        pautas (\texttt{se esperaba un id de clase se encontro ...}).
\end{itemize}

El sintáctico consume tokens bajo demanda: como el símbolo inicial exige llegar al
\texttt{EOF}, todo token del fuente es requerido y, por lo tanto, un error léxico lanzado
por el léxico se propaga a través del sintáctico hasta el módulo principal, que lo reporta
con el mismo formato que un error sintáctico.

\subsection{ExcepcionSintactica}

Excepción (no verificada, al igual que \texttt{ExcepcionLexica}) con el lexema del token
ofensor, su número de línea y un mensaje descriptivo de la naturaleza del error.

\section{Transformación de la gramática a LL(1)}
\label{sec:transformacion}

La gramática de MiniJava no es LL(1) por dos tipos de problemas: producciones con
\textbf{recursión a izquierda} y producciones de un mismo NT cuyas partes derechas
\textbf{comparten primeros} (no factorizadas). Se resolvieron en ese orden: primero la
eliminación de recursión a izquierda y luego la factorización simple y profunda. Las
técnicas aplicadas fueron las vistas en los temas 3.2.3 y 3.2.4. En lo que sigue se
escribe \texttt{eps} por la producción vacía.

\subsection{Eliminación de recursión a izquierda}

\begin{itemize}
  \item \textbf{Expresiones compuestas} (recursión por operador binario a la izquierda):
{\footnotesize
\begin{verbatim}
<ExpresionCompuesta> ::= <ExpresionCompuesta> <OperadorBinario> <ExpresionBasica>
                      | <ExpresionBasica>
\end{verbatim}}
        se transforma en:
{\footnotesize
\begin{verbatim}
<ExpresionCompuesta> ::= <ExpresionBasica> <ECResto>
<ECResto> ::= <OperadorBinario> <ExpresionBasica> <ECResto> | eps
\end{verbatim}}
  \item \textbf{Referencias encadenadas} (recursión a izquierda por \texttt{VarEncadenada},
        \texttt{MetodoEncadenado} y \texttt{AccesoArreglo}):
{\footnotesize
\begin{verbatim}
<Referencia> ::= <Referencia> <VarEncadenada>
               | <Referencia> <MetodoEncadenado>
               | <Referencia> <AccesoArreglo>
               | <Primario>
\end{verbatim}}
        se transforma en:
{\footnotesize
\begin{verbatim}
<Referencia> ::= <Primario> <RResto>
<RResto> ::= . idMetVar <MetodoOVar> <RResto> | <AccesoArreglo> <RResto> | eps
\end{verbatim}}
  \item \textbf{Lista de argumentos formales}:
{\footnotesize
\begin{verbatim}
<ListaArgsFormales> ::= <ListaArgsFormales> , <ArgFormal> | <ArgFormal>
\end{verbatim}}
        se transforma en:
{\footnotesize
\begin{verbatim}
<ListaArgsFormales> ::= <ArgFormal> <LAFResto>
<LAFResto> ::= , <ArgFormal> <LAFResto> | eps
\end{verbatim}}
\end{itemize}

\subsection{Factorización simple y profunda}

\begin{itemize}
  \item \textbf{Miembros de una clase} (\texttt{<Atributo>}, \texttt{<Metodo>} y
        \texttt{<Constructor>} comparten el prefijo \texttt{<Tipo>}; además
        \texttt{<Metodo>} tenía un \texttt{<ModificadorOpcional>} que podía ser
        \texttt{static}). Factorización profunda:
{\footnotesize
\begin{verbatim}
<Miembro> ::= <Tipo> idMetVar <MetodoOAtributo> | <MetodoVoid>
            | static <MetodoResto> | <Constructor>
<MetodoOAtributo> ::= ; | <ArgsFormales> <Bloque>
<MetodoVoid> ::= void idMetVar <ArgsFormales> <Bloque>
<MetodoResto> ::= <MetodoVoid> | <Tipo> idMetVar <ArgsFormales> <Bloque>
\end{verbatim}}
  \item \textbf{Sentencias de expresión}: \texttt{<Asignacion>} y \texttt{<Llamada>} derivan
        ambas en \texttt{<Expresion>}, por lo que se unifican en
        \texttt{<Sentencia> ::= <Expresion> ;}.
  \item \textbf{If con else opcional} (dos producciones con prefijo común):
{\footnotesize
\begin{verbatim}
<If> ::= if ( <Expresion> ) <Sentencia>
       | if ( <Expresion> ) <Sentencia> else <Sentencia>
\end{verbatim}}
        se transforma en:
{\footnotesize
\begin{verbatim}
<If> ::= if ( <Expresion> ) <Sentencia> <ElseOpcional>
<ElseOpcional> ::= else <Sentencia> | eps
\end{verbatim}}
  \item \textbf{Expresión con asignación opcional}:
{\footnotesize
\begin{verbatim}
<Expresion> ::= <ExpresionCompuesta> <OperadorAsignacion> <ExpresionCompuesta>
              | <ExpresionCompuesta>
\end{verbatim}}
        se transforma en:
{\footnotesize
\begin{verbatim}
<Expresion> ::= <ExpresionCompuesta> <ExpresionResto>
<ExpresionResto> ::= <OperadorAsignacion> <ExpresionCompuesta> | eps
\end{verbatim}}
  \item \textbf{Primarios}: \texttt{<AccesoVar>} y \texttt{<LlamadaMetodo>} comparten el
        prefijo \texttt{idMetVar}; \texttt{<CrearArreglo>} y \texttt{<LlamadaConstructor>}
        comparten \texttt{new} y la estructura de \texttt{<TipoBase>} (que a su vez incluye
        \texttt{<TipoReferencia>}):
{\footnotesize
\begin{verbatim}
<Primario> ::= idMetVar <MetodoOVar> | new <ArregloOConstructor> | ...
<MetodoOVar> ::= <ArgsActuales> | eps
<ArregloOConstructor> ::= <TipoReferencia> <AOC2>
                        | <TipoPrimitivo> <DimensionesConTamanio>
                        | idGen <DimensionesConTamanio>
<AOC2> ::= <ArgsActuales> | <DimensionesConTamanio>
\end{verbatim}}
        De la misma forma, en \texttt{<RResto>} los encadenados \texttt{. idMetVar} y
        \texttt{. idMetVar <ArgsActuales>} se factorizan en
        \texttt{. idMetVar <MetodoOVar>}.
  \item \textbf{Dimensiones con tamaño} y \textbf{lista de expresiones} (prefijo común):
{\footnotesize
\begin{verbatim}
<DimensionesConTamanio> ::= [ <Expresion> ] <DCTResto>
<DCTResto> ::= [ <Expresion> ] <DCTResto> | eps
<ListaExps> ::= <Expresion> <ListaExpsResto>
<ListaExpsResto> ::= , <ListaExps> | eps
\end{verbatim}}
\end{itemize}

\subsection{Observaciones sobre la LL(1)-cidad}

\begin{itemize}
  \item \textbf{Else colgante}: es la única construcción residual que, en forma estricta,
        no es LL(1) (Primeros(\texttt{else S}) $\cap$ Siguientes(\texttt{<ElseOpcional>})
        contiene \texttt{else}). El esquema recursivo la resuelve de manera determinista
        consumiendo el \texttt{else} apenas aparece: el \texttt{else} se asocia al
        \texttt{if} más cercano, como en Java.
  \item \textbf{Doble uso de los tokens} \verb!op<! y \verb!op>! (genéricos vs.\ operadores
        relacionales): no genera conflictos porque los genéricos sólo aparecen en contextos
        de tipo (tras \texttt{idClase}) y los siguientes de las reglas de genericidad nunca
        incluyen \verb!op<! ni \verb!op>!. El contexto del parser determina la lectura.
  \item \textbf{Unarios vs.\ binarios} (\texttt{+} y \texttt{-}): los binarios los consume
        \texttt{<ECResto>} después de una \texttt{<ExpresionBasica>} completa, y dentro de
        \texttt{<ExpresionBasica>} los primeros de \texttt{<OperadorUnario>} y de
        \texttt{<Operando>} son disjuntos.
\end{itemize}

\section{Gramática final utilizada}
\label{sec:gramatica}

Versión final LL(1) (equivalente a la original: genera el mismo lenguaje). El símbolo
\texttt{eof} corresponde al token \texttt{EOF} (lexema \verb!$!).

{\footnotesize
\begin{verbatim}
<Inicial> ::= <ListaClases> eof

<ListaClases> ::= <Clase> <ListaClases> | <Interfaz> <ListaClases> | eps

<Clase> ::= class idClase <GenericidadOpcional> <HerenciaOpcional> { <ListaMiembros> }

<Interfaz> ::= interface idClase <GenericidadOpcional> <ExtensionOpcional>
               { <ListaMetodosInterfaz> }

<GenericidadOpcional> ::= < idGen > | eps

<HerenciaOpcional> ::= extends <TipoReferencia> | implements <TipoReferencia> | eps

<ExtensionOpcional> ::= extends <TipoReferencia> | eps

<ListaMiembros> ::= <Miembro> <ListaMiembros> | eps

<ListaMetodosInterfaz> ::= <MetodoInterfaz> <ListaMetodosInterfaz> | eps

<Miembro> ::= <Tipo> idMetVar <MetodoOAtributo>
<Miembro> ::= <MetodoVoid>
<Miembro> ::= static <MetodoResto>
<Miembro> ::= <Constructor>
<MetodoOAtributo> ::= ;
<MetodoOAtributo> ::= <ArgsFormales> <Bloque>
<MetodoVoid> ::= void idMetVar <ArgsFormales> <Bloque>
<MetodoResto> ::= <MetodoVoid>
<MetodoResto> ::= <Tipo> idMetVar <ArgsFormales> <Bloque>

<MetodoInterfaz> ::= <TipoMetodo> idMetVar <ArgsFormales> ;

<Constructor> ::= public idClase <ArgsFormales> <Bloque>

<TipoMetodo> ::= <Tipo> | void

<Tipo> ::= <TipoBase> <DimensionesOpcionales>

<TipoBase> ::= <TipoPrimitivo> | <TipoReferencia> | idGen

<DimensionesOpcionales> ::= [] <DimensionesOpcionales> | eps

<TipoReferencia> ::= idClase <TipoGenericoOpcional>

<TipoPrimitivo> ::= boolean | char | int

<TipoGenericoOpcional> ::= < <InstanciadoOParametrico> > | eps

<InstanciadoOParametrico> ::= idGen | idClase

<ArgsFormales> ::= ( <ListaArgsFormalesOpcional> )

<ListaArgsFormalesOpcional> ::= <ListaArgsFormales> | eps

<ListaArgsFormales> ::= <ArgFormal> <LAFResto>
<LAFResto> ::= , <ArgFormal> <LAFResto> | eps

<ArgFormal> ::= <Tipo> idMetVar

<Bloque> ::= { <ListaSentencias> }

<ListaSentencias> ::= <Sentencia> <ListaSentencias> | eps

<Sentencia> ::= ;
<Sentencia> ::= <Expresion> ;
<Sentencia> ::= <VarLocal> ;
<Sentencia> ::= <Return> ;
<Sentencia> ::= <If>
<Sentencia> ::= <While>
<Sentencia> ::= <Bloque>

<VarLocal> ::= var idMetVar = <ExpresionCompuesta>

<Return> ::= return <ExpresionOpcional>

<ExpresionOpcional> ::= <Expresion> | eps

<If> ::= if ( <Expresion> ) <Sentencia> <ElseOpcional>
<ElseOpcional> ::= else <Sentencia> | eps

<While> ::= while ( <Expresion> ) <Sentencia>

<Expresion> ::= <ExpresionCompuesta> <ExpresionResto>
<ExpresionResto> ::= <OperadorAsignacion> <ExpresionCompuesta> | eps

<OperadorAsignacion> ::= =

<ExpresionCompuesta> ::= <ExpresionBasica> <ECResto>
<ECResto> ::= <OperadorBinario> <ExpresionBasica> <ECResto> | eps

<OperadorBinario> ::= || | && | == | != | < | > | <= | >= | + | - | * | / | %

<ExpresionBasica> ::= <OperadorUnario> <Operando>
<ExpresionBasica> ::= <Operando>

<OperadorUnario> ::= + | - | !

<Operando> ::= <Primitivo>
<Operando> ::= <Referencia>

<Primitivo> ::= true | false | intLiteral | charLiteral | null

<Referencia> ::= <Primario> <RResto>
<RResto> ::= . idMetVar <MetodoOVar> <RResto>
<RResto> ::= <AccesoArreglo> <RResto>
<RResto> ::= eps

<Primario> ::= this
<Primario> ::= stringLiteral
<Primario> ::= idMetVar <MetodoOVar>
<Primario> ::= new <ArregloOConstructor>
<Primario> ::= <LlamadaMetodoEstatico>
<Primario> ::= <ExpresionParentizada>
<ArregloOConstructor> ::= <TipoReferencia> <AOC2>
<ArregloOConstructor> ::= <TipoPrimitivo> <DimensionesConTamanio>
<ArregloOConstructor> ::= idGen <DimensionesConTamanio>
<AOC2> ::= <ArgsActuales> | <DimensionesConTamanio>
<MetodoOVar> ::= <ArgsActuales> | eps

<ExpresionParentizada> ::= ( <Expresion> )

<LlamadaMetodoEstatico> ::= idClase . idMetVar <ArgsActuales>

<DimensionesConTamanio> ::= [ <Expresion> ] <DCTResto>
<DCTResto> ::= [ <Expresion> ] <DCTResto> | eps

<ArgsActuales> ::= ( <ListaExpsOpcional> )

<ListaExpsOpcional> ::= <ListaExps> | eps

<ListaExps> ::= <Expresion> <ListaExpsResto>
<ListaExpsResto> ::= , <ListaExps> | eps

<AccesoArreglo> ::= [ <Expresion> ]
\end{verbatim}
}

\section{Errores sintácticos}

Al detectarse un error se lanza \texttt{ExcepcionSintactica} con el lexema y la línea del
token actual, el módulo principal la captura, reporta el error y finaliza la ejecución (no
hay recuperación: modo pánico fuera del alcance de esta entrega). Los errores léxicos se
reportan con el mismo formato, ya que la excepción \texttt{ExcepcionLexica} se propaga a
través del sintáctico. En el código de error el lexema del token \texttt{EOF} es
\verb!$!.

\begin{verbatim}
Error Sintactico en linea <N>: <mensaje>
[Error:<lexema>|<N>]
\end{verbatim}

El analizador detecta dos tipos de errores sintácticos:

\begin{itemize}
  \item \textbf{Falla de match}: el token actual no coincide con el terminal que exige la
        producción en curso. Mensaje \texttt{se esperaba <esperado> se encontro "<encontrado>"}.
        Por ejemplo, en \verb!class {!:
{\footnotesize
\begin{verbatim}
Error Sintactico en linea 1: se esperaba un id de clase se encontro "{"
[Error:{|1]
\end{verbatim}}
  \item \textbf{Token fuera de los primeros}: el token actual no pertenece a los primeros
        de ninguna producción de un NT sin producción vacía. Mensaje
        \texttt{se encontro "<encontrado>" cuando se esperaba <alternativas>}. Por ejemplo,
        un atributo sin \verb!;!:
{\footnotesize
\begin{verbatim}
Error Sintactico en linea 4: se encontro "}" cuando se esperaba ';' o los
argumentos formales de un metodo
[Error:}|4]
\end{verbatim}}
\end{itemize}

Ejemplos adicionales de reportes:

\begin{center}
{\footnotesize
\setlength{\tabcolsep}{4pt}
\begin{tabular}{|l|l|}
\hline
\textbf{Fuente (fragmento)} & \textbf{Reporte} \\ \hline
\verb!a = b+4 hola;! & \texttt{[Error:hola|5]} (se esperaba \verb!;! tras la expresión) \\ \hline
falta \verb!}! al final del archivo & \verb![Error:$|4]! (EOF inesperado) \\ \hline
\verb!var x;! & \texttt{[Error:;|4]} (se esperaba \verb!=! tras la variable) \\ \hline
\verb!new int[];! & \texttt{[Error:]|4]} (dimensión sin expresión de tamaño) \\ \hline
caracter \verb!#! dentro de un miembro & \verb![Error:#|3]! (error léxico propagado) \\ \hline
\end{tabular}
}
\end{center}

\section{Formato de salida}

\subsection{Análisis exitoso}

\begin{verbatim}
Compilacion Exitosa
[SinErrores]
\end{verbatim}

\subsection{Análisis con errores}

Mensaje descriptivo de la naturaleza del error indicando el lexema del token con el que se
produjo y su número de línea, seguido del código de error (ejemplos en la sección
anterior). Toda la salida se realiza con \texttt{System.out}.

\section{Compilación y ejecución}

No se usa ninguna herramienta de build: solo se necesitan \texttt{javac}, \texttt{jar} y
\texttt{java} de cualquier JDK 21 o superior (probado con JDK 25).

{\footnotesize
\begin{verbatim}
./Compilar.sh             # compila y genera Compilador.jar  (Windows: Compilar.bat)
java -jar Compilador.jar resources/sinErrores/sintCorrecto01.java
./CorrerTests.sh          # compila y ejecuta los testers JUnit  (Windows: CorrerTests.bat)
\end{verbatim}
}

\texttt{CorrerTests.sh} descarga automáticamente \texttt{junit-4.13.2.jar} y
\texttt{hamcrest-core-1.3.jar} en \texttt{lib/} si no están (requiere \texttt{curl});
si no hay red, se pueden descargar a mano de Maven Central y dejarlos en esa carpeta.

\subsection{Comandos manuales equivalentes}

Compilar (desde \texttt{codigo/}):

{\footnotesize
\begin{verbatim}
javac -d build/classes \
  src/main/java/moduloprincipal/*.java \
  src/main/java/analizadorlexico/*.java \
  src/main/java/analizadorsintactico/*.java \
  src/main/java/sourcemanager/*.java
jar cfe Compilador.jar moduloprincipal.ModuloPrincipal -C build/classes .
\end{verbatim}
}

Correr los tests (desde \texttt{codigo/}, con JUnit 4.13.2 y Hamcrest 1.3 en \texttt{lib/}):

{\footnotesize
\begin{verbatim}
javac -cp "build/classes:lib/junit-4.13.2.jar:lib/hamcrest-core-1.3.jar" \
  -d build/test-classes src/test/java/test/*.java

java -cp "build/classes:build/test-classes:lib/junit-4.13.2.jar:\
lib/hamcrest-core-1.3.jar" org.junit.runner.JUnitCore \
  test.TesterDeCasosSinErrores test.TesterDeCasosConErrores
\end{verbatim}
}

En Windows el separador del classpath es \texttt{;} en lugar de \texttt{:}.

Los testers usan el working directory \texttt{codigo/} para acceder a
\texttt{resources/sinErrores/} y \texttt{resources/conErrores/} (los scripts ya se
posicionan ahí; los comandos manuales deben correrse desde \texttt{codigo/}). Los casos
sin errores deben ser aceptados (la salida debe contener \texttt{[SinErrores]}); los casos
con errores llevan en su primera línea el código esperado con el formato
\verb!///[Error:lexema|nroLinea]! y el tester verifica que aparezca en la salida.

\section{Tests}

Se incluyen los casos provistos por la cátedra y casos propios desarrollados para esta
etapa (el diseño de los casos de prueba es responsabilidad del alumno). Resultado actual:
\textbf{26/26 casos OK} (11 sin errores y 15 con errores).

\subsection{Casos provistos por la cátedra}

\begin{itemize}
  \item \texttt{sintCorrecto01.java} - clase vacía.
  \item \texttt{sintCorrecto02.java} - método dinámico sin argumentos y cuerpo vacío.
  \item \texttt{sintCorrecto03.java} - método estático con parámetro y bloque anidado con
        sentencia vacía.
  \item \texttt{sintCorrecto04.java} - asignación con \texttt{this} encadenado, atributo y
        constructor.
  \item \texttt{sintError01.java} - herencia doble (\verb!,! tras el extends).
  \item \texttt{sintError02.java} - id de clase en minúsculas.
  \item \texttt{sintError03.java} - expresión mal formada (falta operador).
  \item \texttt{sintError04.java} - \texttt{else} en lugar de sentencia.
\end{itemize}

\subsection{Casos propios}

\begin{itemize}
  \item \texttt{sintCorrecto05.java} - interfaz con genericidad y extensión, varios
        métodos.
  \item \texttt{sintCorrecto06.java} - clase genérica con herencia genérica, atributos de
        todos los tipos (primitivo, arreglo, clase genérica, parámetro de tipo) y
        constructor.
  \item \texttt{sintCorrecto07.java} - if/else anidados, while, var, return con y sin
        expresión, bloques anidados y sentencia vacía.
  \item \texttt{sintCorrecto08.java} - expresiones unarias, binarias y parentizadas con
        todos los operadores y literales.
  \item \texttt{sintCorrecto09.java} - referencias encadenadas, creación de arreglos con
        \texttt{new}, accesos con índice y \texttt{this}.
  \item \texttt{sintCorrecto10.java} - constructores, \texttt{new} de clases con y sin
        genericidad, \texttt{new} de parámetro de tipo y llamadas a métodos estáticos.
  \item \texttt{sintCorrecto11.java} - \texttt{implements}, tipos arreglo en parámetros y
        retorno, varios argumentos formales y actuales.
  \item \texttt{sintError05.java} - falta el id de clase (\verb!class {!).
  \item \texttt{sintError06.java} - EOF inesperado (falta la llave de cierre), reporta
        lexema \verb!$!.
  \item \texttt{sintError07.java} - falta \verb!;! en un atributo.
  \item \texttt{sintError08.java} - operador binario sin operando derecho.
  \item \texttt{sintError09.java} - método de interfaz con cuerpo (se esperaba \verb!;!).
  \item \texttt{sintError10.java} - \verb!new int[]! sin tamaño de dimensión.
  \item \texttt{sintError11.java} - paréntesis del \texttt{if} sin cerrar.
  \item \texttt{sintError12.java} - \texttt{var} sin inicializador.
  \item \texttt{sintError13.java} - \texttt{extends} sin tipo.
  \item \texttt{sintError14.java} - acceso encadenado sin identificador (\verb!a. = 3;!).
  \item \texttt{sintError15.java} - error léxico detectado durante el análisis
        sintáctico (propagación de \texttt{ExcepcionLexica}).
\end{itemize}

\section{Logros}

Para esta etapa se implementa el analizador sintáctico base sobre la gramática LL(1) de
MiniJava; \textbf{no se intentan logros opcionales de la Etapa 2} (expresiones lambda,
variables locales clásicas, visibilidad explícita, recuperación en modo pánico,
\texttt{for}, genericidad avanzada/anidada, atributos inicializados, inicialización de
arreglos, operador ternario ni operadores posfijos).

De la etapa anterior se conserva vigente el logro del \textbf{Manejador de Archivos
Eficiente} (\texttt{SourceManagerEficiente}), que sigue siendo el módulo de lectura del
fuente. La multi-detección de errores léxicos de la Etapa 1 no aplica en esta etapa: el
módulo principal ahora finaliza al primer error detectado (sintáctico o léxico), según lo
exigen las pautas de la Etapa 2.

\end{document}
